\section{Auditoría de fuentes y pregunta central del curso}

En apariencia, esta clase va del póquer y el go hasta Diplomacy, un juego de mesa para siete jugadores; lo que en realidad recorre toda la sesión son dos problemas de sistemas. Primero, una vez terminado el entrenamiento, cuánto cómputo debe invertir un modelo en una sola decisión. Segundo, cuando una tarea requiere a la vez modelar el mundo, optimizar una estrategia y comunicarse en lenguaje natural, qué responsabilidades deben asignarse a un mismo modelo y cuáles deben separarse necesariamente. Noam Brown usa la búsqueda para mostrar que \term{el presupuesto de cómputo debe variar según el valor de la decisión}, y luego usa Cicero para mostrar que \term{la fluidez lingüística no sustituye el ciclo cerrado estratégico}.

\subsection{Cómo se reconstruye esta clase}

El sitio oficial del curso no ofrece un slide deck independiente para descargar, por lo que esta clase no puede tomar como fuente visual borradores anteriores ni resúmenes de segunda mano publicados en internet. Tomamos como original maestro la grabación oficial en 1080p de Stanford Online, muestreamos la región de screen-share cada dos segundos, de 1,600 cuadros seleccionamos 70 candidatos con alta exhaustividad y, tras eliminar duplicados manualmente, quedaron 42 imágenes con valor didáctico propio. Las explicaciones orales del texto provienen de los subtítulos oficiales en inglés elaborados por personas, con un total de 1,176 intervalos de tiempo; las fórmulas y los términos se cotejaron con los artículos primarios recomendados en clase.

\begin{longtable}{L{0.22\textwidth}L{0.31\textwidth}L{0.37\textwidth}}
\toprule
Capa de evidencia & Material local & Función en los apuntes \\
\midrule
Grabación oficial & Video de 53:20 y 42 diapositivas recuperadas & Determina el orden de la clase, las gráficas, la arquitectura del sistema y los casos de demostración \\
Subtítulos oficiales elaborados por personas & \texttt{lecture14.en.srt} & Recuperan la motivación, las analogías de clase, las restricciones, la sesión de Q\&A y el tono del ponente \\
Artículos recomendados por el curso & Cicero, piKL, No-Press Diplomacy & Permiten cotejar los objetivos de los algoritmos, los límites de los experimentos y el significado de los términos \\
Archivos locales de auditoría & manifest, coverage, teacher-voice ledger & Demuestran que cada source node tiene un destino explícito en el texto \\
\bottomrule
\end{longtable}

\begin{warningbox}{Límite histórico}
Esta clase reconstruye el contenido impartido hasta el 31 de enero de 2023. Los reasoning models, las test-time scaling laws o los productos de agent de propósito general posteriores no pueden proyectarse hacia atrás como hechos de la clase. Al final se incluye por separado una sección de «Ampliación a sistemas modernos», claramente separada de las afirmaciones que hizo el ponente en ese momento.
\end{warningbox}

\subsection{Cuatro preguntas que recorren toda la clase}

Al leer más adelante las curvas de póquer y la arquitectura de Cicero, no conviene quedarse en el nivel de «la IA vuelve a vencer a los humanos». Vale más seguir estas preguntas: si la calidad de la decisión proviene de la base policy entrenada o de la search realizada en el momento; si el comportamiento humano es el objetivo de optimización, una distribución a priori o un modelo del entorno; si el modelo de lenguaje decide la estrategia o solo expresa una estrategia ya decidida; y cómo rechaza el sistema un mensaje gramaticalmente correcto pero que perjudicaría la recompensa futura.

\begin{importantbox}{La tesis sistémica de esta clase}
Un agente inteligente sólido no es «un modelo de lenguaje más grande con un prompt», sino un ciclo cerrado sujeto a un presupuesto: predice a los demás participantes del entorno, planifica su propia estrategia, condensa esa estrategia en un intent comunicable, genera mensajes candidatos, evalúa cómo cambiarían esos mensajes el comportamiento de los demás y vuelve a planificar.
\end{importantbox}

\subsection{Glosario: distinguir primero los conceptos afines}

Esta clase alterna una y otra vez entre cuatro niveles: entrenamiento, inferencia, estrategia y lenguaje. La tabla siguiente reúne primero los términos clave, para no confundir «un modelo más grande», «una búsqueda más profunda», «más parecido a un humano» y «mejor comunicación» como si fueran la misma capacidad.

\begin{longtable}{L{0.23\textwidth}L{0.30\textwidth}L{0.37\textwidth}}
\toprule
Término & Problema que resuelve & Mecanismo central y relación con el curso \\
\midrule
Base policy & Cómo producir rápidamente una distribución de acciones sin búsqueda & Se obtiene mediante entrenamiento; es la distribución a priori de candidatos para la search, no la respuesta final \\
Inference-time search & Si la decisión actual justifica gastar más cómputo & Con los pesos fijos, despliega, evalúa y compara candidatos; el presupuesto puede variar según el valor \\
Exploitability & Cuánto puede aprovechar una estrategia un oponente óptimo & Mide la distancia en el peor caso respecto al Nash equilibrium; es más robusta que la tasa de victorias común \\
Human anchor policy & Cómo conservar convention de acción comprensibles para las personas & Imita datos humanos; aporta una distribución a priori de coordinación, pero también hereda las debilidades humanas \\
Intent & Qué plan debe expresar exactamente el modelo de lenguaje & Cicero usa actions ejecutables a corto plazo como condición estructurada de generación \\
Grounded dialogue & Cómo se vinculan los mensajes con el estado compartido del mundo & Cada mensaje debe corresponder al board, la history, el plan y las posibles consecuencias en el comportamiento \\
Value-based filtering & Cómo rechazar mensajes fluidos pero estratégicamente dañinos & Predice el efecto del mensaje en la policy de los demás y luego evalúa el expected value del plan \\
\bottomrule
\end{longtable}

\subsection{Resumen de la sección}

Esta clase se apoya en tres capas de evidencia: la grabación oficial, los subtítulos elaborados por personas y los artículos primarios. A continuación se explica primero «por qué vale la pena gastar en search durante la inferencia» y después se aborda Diplomacy, donde la comunicación, la cooperación y la competencia coexisten.
